% Canonical Paper I source.
Arithmetic evaluation is usually presented as a map from a syntactic expression
to its value.  This description suppresses the order in which intermediate
values are formed.  We begin by isolating the finite binary trees for which the
internal operations have only one legal temporal order.  Such a tree is not yet
a history: one must also mark which leaf of its innermost operation is to be
carried forward as the evolving value.  The resulting object, a marked spinal
history, is the syntactic unit used in the remainder of the paper.

Throughout this section, $K$ is a field and
\[
  \Omega=\{+,-,\times,/\}
\]
is the binary arithmetic signature.  Division is interpreted as a partial
operation in ordinary arithmetic.  The combinatorial statements below apply
to any binary signature; the choice of $\Omega$ is needed only when histories
are evaluated.

\subsection{Expression trees and dependency}

\begin{definition}[Arithmetic expression tree]
\label{def:arithmetic-expression-tree}
An \emph{arithmetic expression tree over $K$} is a finite rooted planar binary
tree whose leaves are labelled by elements of $K$ or by formal variables and
whose internal vertices are labelled by elements of $\Omega$.  Its set of
internal vertices is denoted by $I(T)$.
\end{definition}

The planar structure records the first and second operand slots.  It does not,
by itself, specify a temporal order between operations on disjoint subtrees.
That order is recorded by the dependency relation.

\begin{definition}[Internal dependency poset]
\label{def:internal-dependency-poset}
For $u,v\in I(T)$, write $u\preccurlyeq_T v$ when either $u=v$ or the value
produced at $u$ is required as an input at $v$.  Equivalently, $v$ lies on the
path from $u$ to the root.  A \emph{legal internal evaluation order} is a
linear extension of the finite poset $(I(T),\preccurlyeq_T)$.
\end{definition}

Leaves are not included in $I(T)$.  Thus a tree with no internal vertex has the
unique empty internal evaluation order.  This convention separates the order
of operations from the prior availability of atomic inputs.

We shall use the following elementary finite-poset fact in both directions.

\begin{lemma}[Unique linear extensions]
\label{lem:unique-linear-extension}
A finite poset has a unique linear extension if and only if it is a chain.
\end{lemma}

\begin{proof}
Every linear extension of a chain is its given total order, so a chain has a
unique linear extension.  Conversely, let $P$ be a finite poset and suppose
that $u,v\in P$ are incomparable.  Adding the relation $u<v$ and taking its
transitive closure cannot create a directed cycle: such a cycle would contain
a pre-existing path $v\leq u$, contrary to incomparability.  The resulting
finite partial order has a linear extension.  The same argument with $v<u$
gives a second linear extension.  One orders $u$ before $v$ and the other
orders $v$ before $u$, so the extension of $P$ is not unique.  Hence a poset
with a unique linear extension contains no incomparable pair and is a chain.
\end{proof}

\begin{theorem}[Sequential-tree classification]
\label{thm:sequential-tree-classification}
For a finite binary expression tree $T$, the following conditions are
equivalent.
\begin{enumerate}
  \item The dependency poset $I(T)$ has a unique linear extension.
  \item The dependency poset $I(T)$ is a chain.
  \item Every internal vertex has at most one internal child.
  \item The internal vertices form a single spine: either $I(T)$ is empty, or
  there is a path from the root to an innermost internal vertex containing
  every internal vertex of $T$.
\end{enumerate}
\end{theorem}

\begin{proof}
The equivalence of the first two conditions is
Lemma~\ref{lem:unique-linear-extension}.

Suppose that $I(T)$ is a chain.  If an internal vertex had two internal
children, neither child could depend on the other, since they lie in disjoint
subtrees.  They would therefore be incomparable in $I(T)$, a contradiction.
Thus (2) implies (3).

Assume (3).  If $I(T)$ is empty there is nothing to prove.  Otherwise begin at
the root.  Whenever the current internal vertex has an internal child, that
child is unique.  Following these children produces a path ending at an
internal vertex with two leaf children.  Every internal vertex lies on this
path: an internal vertex off the path would occur in the other child-subtree
of some vertex on the path, giving that vertex two internal children.  Hence
the path contains all internal vertices, proving (4).

Finally, if all internal vertices lie on one root-to-leaf path, their ancestor
relation is total.  Read from the innermost vertex towards the root, each
operation depends on its predecessor.  Thus $I(T)$ is a chain, proving
(4)$\Rightarrow$(2) and completing the equivalence.
\end{proof}

Figure~\ref{fig:sequential-branching-dependency} contrasts the two possible
local patterns at the level of internal dependencies.  It is only a visual
summary of the order relation; the classification rests on the preceding
proof.

\begin{figure}[t]
\centering
\begin{tikzpicture}[
  scale=0.92,
  every node/.style={font=\small},
  iv/.style={circle,draw=black!70,fill=blue!7,minimum size=7mm,inner sep=0pt},
  dep/.style={-{Latex[length=2mm]},thick,black!65}
]
  \node[font=\small\bfseries] at (-2.4,1.45) {Sequential spine};
  \node[iv] (s3) at (-2.4,0.85) {$v_3$};
  \node[iv] (s2) at (-2.4,0) {$v_2$};
  \node[iv] (s1) at (-2.4,-0.85) {$v_1$};
  \draw[dep] (s1) -- (s2);
  \draw[dep] (s2) -- (s3);
  \node[below=3pt of s1,font=\scriptsize] {first};

  \node[font=\small\bfseries] at (2.4,1.45) {Branching dependency};
  \node[iv] (b3) at (2.4,0.65) {$u_3$};
  \node[iv] (b1) at (1.65,-0.55) {$u_1$};
  \node[iv] (b2) at (3.15,-0.55) {$u_2$};
  \draw[dep] (b1) -- (b3);
  \draw[dep] (b2) -- (b3);
  \node[anchor=north east,font=\scriptsize] at (1.45,-0.92) {incomparable};
  \node[anchor=north west,font=\scriptsize] at (3.35,-1.18) {incomparable};
\end{tikzpicture}
\caption{A sequential tree has a chain of internal dependencies (left).  A
branching tree contains incomparable internal operations and therefore admits
more than one legal internal evaluation order (right).}
\label{fig:sequential-branching-dependency}
\end{figure}

\begin{definition}[Sequential tree]
\label{def:sequential-tree}
A finite binary expression tree satisfying the equivalent conditions of
Theorem~\ref{thm:sequential-tree-classification} is called a
\emph{sequential tree}.
\end{definition}

The adjective refers only to internal evaluation.  It does not choose an
active leaf, and it does not assert that ordinary arithmetic evaluation is
defined.  For example,
\[
  (a+b)(c+d)
\]
is not sequential: the two additions are incomparable dependencies of the
final multiplication.  By contrast,
\[
  \bigl((a+b)c-d\bigr)/e
\]
is sequential, because its internal vertices lie on one spine.

\subsection{One-hole contexts and the marked accumulator}

Let $z$ denote a hole rather than an arithmetic variable.  For
$\omega\in\Omega$ and $c\in K$, define the two elementary one-hole contexts
\begin{equation}
  C_{\omega,c}^{(1)}[z]=\omega(z,c),
  \qquad
  C_{\omega,c}^{(2)}[z]=\omega(c,z).
  \label{eq:elementary-one-hole-contexts}
\end{equation}
The superscript gives the operand slot occupied by the evolving value.  It is
therefore intrinsic to the marked process and should not be replaced by a
description based on how a planar drawing happens to expand.

At the innermost internal vertex of a nontrivial sequential tree, both children
are leaves.  The dependency poset determines that this operation is performed
first, but it does not determine which of its two inputs is the value to be
carried along the spine.  Choosing one of those leaves as the initial
accumulator removes this ambiguity.

\begin{definition}[Marked spinal history]
\label{def:marked-spinal-history}
A \emph{free marked spinal history over $K$} is a finite chronological word
\[
  \gamma=(g_1,\ldots,g_n),
  \qquad
  g_i=C_{\omega_i,c_i}^{(\varepsilon_i)},
  \quad \varepsilon_i\in\{1,2\}.
\]
The context $g_1$ is applied first and $g_n$ last.  Its \emph{chirality word}
is
\[
  \varepsilon(\gamma)=\varepsilon_1\cdots\varepsilon_n.
\]
A \emph{bounded marked spinal history} is a pair $(x_0;\gamma)$ consisting of
an initial accumulator $x_0$ and a free history.  The empty word is allowed;
its bounded form represents the marked single-leaf tree.
\end{definition}

When no confusion is possible, we call either form a marked spinal history and
state explicitly whether an initial value is present.  Earlier drafts called
certain pure-chirality cases threadlike expressions.  The present definition
also includes mixed chirality and makes the marking part of the data.

Figure~\ref{fig:slot-chirality-histories} displays operand-slot chirality
directly in the one-hole contexts.  The diagram is a notation guide, not a
substitute for the marked-tree correspondence proved below.

\begin{figure}[t]
\centering
\begin{tikzpicture}[
  box/.style={draw=black!65,rounded corners=2pt,fill=blue!5,
              minimum height=8mm,inner xsep=6pt,font=\small},
  chronarrow/.style={-{Latex[length=2mm]},thick,black!60},
  rowlabel/.style={font=\small\bfseries,anchor=east}
]
  \node[rowlabel] at (-3.6,1.25) {slot-(1)};
  \node[box] (a11) at (-1.9,1.25) {$\omega_1(\bullet,c_1)$};
  \node[box] (a12) at (0.55,1.25) {$\omega_2(\bullet,c_2)$};
  \draw[chronarrow] (a11) -- (a12);
  \node[anchor=west,font=\small] at (2.05,1.25) {$\varepsilon=11$};

  \node[rowlabel] at (-3.6,0) {slot-(2)};
  \node[box] (a21) at (-1.9,0) {$\omega_1(c_1,\bullet)$};
  \node[box] (a22) at (0.55,0) {$\omega_2(c_2,\bullet)$};
  \draw[chronarrow] (a21) -- (a22);
  \node[anchor=west,font=\small] at (2.05,0) {$\varepsilon=22$};

  \node[rowlabel] at (-3.6,-1.25) {mixed};
  \node[box] (am1) at (-2.35,-1.25) {$\omega_1(\bullet,c_1)$};
  \node[box] (am2) at (0,-1.25) {$\omega_2(c_2,\bullet)$};
  \node[box] (am3) at (2.35,-1.25) {$\omega_3(\bullet,c_3)$};
  \draw[chronarrow] (am1) -- (am2);
  \draw[chronarrow] (am2) -- (am3);
  \node[anchor=west,font=\small] at (3.25,-1.25) {$\varepsilon=121$};
\end{tikzpicture}
\caption{Pure slot-(1), pure slot-(2), and mixed marked histories.  In each
box, $\bullet$ is the current accumulator---the marked seed in the first box
and the preceding output thereafter.  Its displayed position is the operand
slot, while time always runs along the arrows.}
\label{fig:slot-chirality-histories}
\end{figure}

\begin{proposition}[Marked-tree--history correspondence]
\label{prop:marked-tree-history-correspondence}
Let $T$ be a sequential tree all of whose leaves are labelled by elements of
$K$.
\begin{enumerate}
  \item If $T$ is a single leaf, marking that leaf determines the empty bounded
  history.
  \item If $T$ is nontrivial, a choice of one leaf at its innermost internal
  vertex as the initial accumulator determines a unique bounded marked spinal
  history.
  \item Conversely, a bounded marked spinal history determines a unique marked
  planar sequential tree.
\end{enumerate}
These two constructions are mutually inverse when operation labels, constants,
operand slots, and the initial mark are retained.
\end{proposition}

\begin{proof}
The single-leaf case is immediate.  Suppose $T$ is nontrivial.  By
Theorem~\ref{thm:sequential-tree-classification}, its internal vertices form a
unique spine.  At the innermost vertex, mark one leaf $x_0$ and denote the
other by $c_1$.  The operation label and the slot containing $x_0$ determine
$g_1=C_{\omega_1,c_1}^{(\varepsilon_1)}$.  Move one vertex towards the root.
Its off-spine child must be a leaf, since otherwise it would be a second
internal child.  Its label $c_2$, the operation label, and the slot occupied by
the spine determine $g_2$.  Repeating this procedure reaches the root and
produces a unique chronological word $(g_1,\ldots,g_n)$.

Conversely, start with a marked leaf $x_0$.  Replace it by the planar binary
tree $C_{\omega_1,c_1}^{(\varepsilon_1)}[x_0]$, then place the resulting tree
in slot $\varepsilon_2$ of the next context, and continue.  Exactly one child
at every new internal vertex is internal, so the result is sequential.  Each
step records the operation, off-spine leaf, and slot used in the forward
construction.  The constructions therefore recover one another.
\end{proof}

The proposition concerns fully labelled planar trees.  If parentheses,
constant labels, or the marked seed are forgotten, the correspondence is no
longer injective.

\subsection{Composition and ordinary evaluation}

If
\[
  \gamma=(g_1,\ldots,g_m),
  \qquad
  \delta=(h_1,\ldots,h_n),
\]
their chronological concatenation is
\begin{equation}
  \gamma\delta=(g_1,\ldots,g_m,h_1,\ldots,h_n),
  \label{eq:chronological-history-concatenation}
\end{equation}
so $\gamma$ is performed first and $\delta$ second.  This convention is fixed
for the entire paper and is summarized again in Appendix~\ref{app:conventions}.

Each elementary context determines a partial function $g_i\colon K\dashrightarrow
K$.  Starting from $x\in K$, define recursively
\[
  x_0=x,
  \qquad
  x_i=g_i(x_{i-1}).
\]
The bounded history is \emph{ordinarily admissible at $x$} when every $x_i$
is defined in ordinary field arithmetic.  In that case its endpoint is
\begin{equation}
  \nu_x(\gamma)
  =g_n\circ\cdots\circ g_1(x).
  \label{eq:ordinary-history-evaluation}
\end{equation}
The admissible domain of a free history is the set of such $x$.  In particular,
a second-slot division $c/z$ excludes any state at which the evolving value is
zero, while a first-slot division $z/c$ requires $c\neq0$.

The concatenation convention gives
\begin{equation}
  \nu_x(\gamma\delta)
  =\nu_{\nu_x(\gamma)}(\delta)
  \label{eq:ordinary-evaluation-concatenation}
\end{equation}
whenever both sides are ordinarily admissible.  Equation
\eqref{eq:ordinary-evaluation-concatenation} is a statement about partial
functions.  Chapter~\ref{sec:projective-affine} introduces a projective
evaluation with a different domain convention.

\subsection{Mirror, temporal reversal, and path inverse}

Three operations on a context word must be kept separate.

\begin{definition}[Mirror and temporal reversal]
\label{def:mirror-temporal-reversal}
For $g=C_{\omega,c}^{(\varepsilon)}$, let
\[
  m(g)=C_{\omega,c}^{(3-\varepsilon)}.
\]
The \emph{mirror} and \emph{temporal reversal} of
$\gamma=(g_1,\ldots,g_n)$ are respectively
\[
  m\gamma=(m(g_1),\ldots,m(g_n)),
  \qquad
  r\gamma=(g_n,\ldots,g_1).
\]
The mirror exchanges operand slots while retaining chronological order;
temporal reversal reverses chronological order while retaining the contexts.
\end{definition}

Both operations are involutions on marked words, and they commute:
$m(r\gamma)=r(m\gamma)$.  This is a syntactic statement; ordinary domains of
the evaluated words need not agree.

\begin{definition}[Path inverse]
\label{def:path-inverse}
Suppose every context $g_i$ in $\gamma$ acts as an invertible function on a
declared domain.  The \emph{path inverse} is the functional word
\[
  \gamma^{-1}=(g_n^{-1},\ldots,g_1^{-1}).
\]
It is defined only under this invertibility hypothesis and need not have the
same operation labels or ordinary admissible domain as $r\gamma$.
\end{definition}

For example, let $g_1[z]=z+1$ and $g_2[z]=2z$.  Then
\[
  \gamma=(g_1,g_2)\colon x\longmapsto2x+2.
\]
Because addition and multiplication are symmetric operations,
$m\gamma$ has the displayed expression $2(1+x)$ and induces the same operator
$2x+2$.  In contrast,
\[
  r\gamma=(g_2,g_1)\colon x\longmapsto2x+1,
\]
while
\[
  \gamma^{-1}=(z\mapsto z/2,\ z\mapsto z-1)
  \colon x\longmapsto x/2-1.
\]
Thus equality with a mirror in this example is caused by the symmetry of the
chosen elementary operations; it does not identify mirror, reversal, and
inverse.  The affine order defect studied later compares temporal orders.  It
is not automatically an invariant of planar mirror.

\subsection{Information retained and information forgotten}

Several equivalence relations occur naturally, and the remainder of the paper
uses them at different stages.  It is useful to order them by the information
they forget.

\begin{enumerate}
  \item \emph{Literal syntactic equality} retains the fully parenthesized
  written expression.
  \item \emph{Marked-history equality} retains the initial mark, the
  chronological context word, and every operand-slot label.  It may forget
  inessential typography but not process data.
  \item \emph{Operator equality} identifies histories inducing the same
  partial function on a stated ordinary domain, or the same projective
  transformation after projective evaluation.
  \item \emph{Charge equality}, when a charge map has been specified, retains
  only the chosen abelian totals of a history.
  \item \emph{Endpoint equality} at $x$ asks only whether the admissible values
  $\nu_x(\gamma)$ and $\nu_x(\delta)$ coincide.
  \item \emph{Quotient or relational equality} is any further identification
  imposed by an explicitly declared rewriting or condensation relation.
\end{enumerate}

This list is not a claim that every adjacent pair gives a strict quotient in
all languages.  It is a discipline for stating which level is being used.
For example, the two histories
\[
  \gamma=(z\mapsto z+1,\ z\mapsto2z),
  \qquad
  \delta=(z\mapsto3z,\ z\mapsto z+1)
\]
are distinct marked histories and induce the distinct operators $2z+2$ and
$3z+1$, but at the initial value $x=1$ both have endpoint $4$.

\subsection{Representative chirality patterns}

The slot notation exposes examples that the older comb terminology obscured.

\begin{example}[Pure slot-(1) history]
The word
\[
  C_{+,2}^{(1)},\quad C_{\times,3}^{(1)},\quad C_{-,4}^{(1)}
\]
has chirality word $111$ and evaluates ordinarily as
$x\mapsto (x+2)3-4$.
\end{example}

\begin{example}[Pure slot-(2) history]
The word
\[
  C_{-,2}^{(2)},\quad C_{/,3}^{(2)}
\]
has chirality word $22$ and evaluates as
\[
  x\longmapsto \frac{3}{2-x},
\]
on the ordinary domain $K\setminus\{2\}$.  The accumulator occupies the
second operand slot at both steps.
\end{example}

\begin{example}[Mixed chirality]
The word
\[
  C_{-,1}^{(1)},\quad C_{/,2}^{(2)},\quad C_{+,3}^{(1)}
\]
has chirality word $121$ and evaluates as
\[
  x\longmapsto \frac{2}{x-1}+3,
\]
on its ordinary domain.  Its dependency poset is still a chain; mixed
chirality does not introduce branching.
\end{example}

The syntactic passage is therefore
\[
  \begin{gathered}
    \text{unmarked sequential tree}
    \xrightarrow{\text{choice of accumulator}}
    \text{marked planar sequential tree},\\
    \text{marked planar sequential tree}
    \longleftrightarrow
    \text{bounded marked spinal history},\\
    \text{bounded marked spinal history}
    \longrightarrow
    \text{operator}
    \longrightarrow
    \text{endpoint}.
  \end{gathered}
\]
when the history is bounded and all planar labels are retained.  The first arrow is a
choice and the double arrow is a bijection; only the later evaluation arrows forget
process data.  Chapter~\ref{sec:projective-affine}
constructs the operator stage for all non-degenerate bilateral contexts.
